% 图契约：圆柱电阻表面的能流；几何、方向和公式与正文同一约定。
\documentclass[tikz,border=5pt]{standalone}
\usepackage{fontspec,xeCJK,amsmath,amssymb}
\IfFontExistsTF{Times New Roman}{\setmainfont{Times New Roman}}{\setmainfont{TeX Gyre Termes}}
\IfFontExistsTF{Heiti SC}{\setCJKmainfont[AutoFakeBold=1.5]{Heiti SC}}{\setCJKmainfont[AutoFakeBold=1.5]{FandolHei-Regular}}
\usetikzlibrary{arrows.meta,calc,patterns}
\definecolor{ink}{HTML}{243745}
\definecolor{blue}{HTML}{275DAB}
\definecolor{green}{HTML}{177C63}
\definecolor{red}{HTML}{BD4935}
\definecolor{orange}{HTML}{AA6C16}
\definecolor{purple}{HTML}{715393}
\begin{document}
\begin{tikzpicture}[x=1cm,y=1cm,font=\fontsize{9.5}{12}\selectfont,
v/.style={-{Latex[length=2mm,width=1.3mm]},line width=.85pt},
axis/.style={-{Latex[length=1.7mm,width=1.1mm]},ink!65,line width=.55pt},
aux/.style={ink!45,dashed,line width=.55pt},
lab/.style={fill=white,inner sep=1.8pt},
card/.style={draw=ink!25,fill=blue!3,rounded corners=2pt,text width=6.1cm,minimum height=.85cm,inner sep=6pt,align=center}]
\path[use as bounding box](-8,-3.6)rectangle(8,3.65);
\node[font=\bfseries\fontsize{10.5}{13}\selectfont,text width=15.1cm,align=center]at(0.00000,3.22000){M7-006　圆柱电阻表面的能流};
\begin{scope}[shift={(-4,0)}]
\fill[blue!5](-2,-1)rectangle(2,1);\draw[ink](-2,-1)rectangle(2,1);
\draw[v,blue](-2.60000,0.00000)--(2.80000,0.00000);
\node[]at(2.50000,-0.50000){$+z,I$};
\node[lab,text=blue]at(0.00000,0.40000){$\mathbf E,\mathbf J$};
\draw[v,red](-1.40000,1.90000)--(-1.40000,1.05000);
\draw[v,red](-1.40000,-1.90000)--(-1.40000,-1.05000);
\draw[v,red](0.00000,1.90000)--(0.00000,1.05000);
\draw[v,red](0.00000,-1.90000)--(0.00000,-1.05000);
\draw[v,red](1.40000,1.90000)--(1.40000,1.05000);
\draw[v,red](1.40000,-1.90000)--(1.40000,-1.05000);
\draw[green](0.7,1.4)circle(0.13);
\fill[green](0.7,1.4)circle(.035);
\node[anchor=west,text=green]at(0.92000,1.40000){$\mathbf H$};
\node[text=red]at(-1.75000,2.25000){$\mathbf S$ 径向流入};
\draw[{Latex[length=1.3mm]}-{Latex[length=1.3mm]},ink!65,line width=.5pt](-2.00000,-2.35000)--(2.00000,-2.35000);
\node[lab]at(0.00000,-2.60000){$\ell$};
\draw[{Latex[length=1.3mm]}-{Latex[length=1.3mm]},ink!65,line width=.5pt](2.45000,0.00000)--(2.45000,1.00000);
\node[lab]at(2.80000,0.50000){$a$};
\end{scope}
\node[card]at(4.00000,1.25000){$E_z=I/(\sigma\pi a^2)$};
\node[card]at(4.00000,-1.05000){$\mathbf S(a)=-\frac{I^2}{2\sigma\pi^2a^3}\mathbf e_r$};
\node[text width=15.2cm,align=center]at(0.00000,-3.20000){对侧表面 $2\pi a\ell$ 积分，得到 $P_{\rm in}=I^2R$。};
\end{tikzpicture}\end{document}
